\section*{Chapter \ref{ch:g11:unemployment} --- Unemployment}
% figdata: figdata/grade-11/12-unemployment/out/ (eu_unemp, jst_*, bren_path, bren_key); Brenland numbers invented.

\begin{solution}{exo:g11:unemployment:1}
Inactive; unemployed; employed; inactive (a discouraged worker: he wants a job but is not looking).
\end{solution}

\begin{solution}{exo:g11:unemployment:2}
Unemployed (found a job that starts within three months, and available); inactive; employed, and underemployed; inactive (not available within two weeks).
\end{solution}

\begin{solution}{exo:g11:unemployment:3}
Labour force $600 + 40 = 640$; unemployment rate $40/640 = 6.25\,\%$; \omterm{def:g11:unemployment:employment-rate}{employment rate} $600/1000 = 60\,\%$.
\end{solution}

\begin{solution}{exo:g11:unemployment:4}
160 unemployed and 1840 employed; \omterm{def:g11:unemployment:employment-rate}{employment rate} $1840/3000 \approx 61.3\,\%$.
\end{solution}

\begin{solution}{exo:g11:unemployment:5}
Highest: Spain, 26.5\,\%; lowest: Germany, 6.6\,\%. Germany: $6.6/3.5 \approx 1.9$; Sweden: $24.3/8.4 \approx 2.9$.
\end{solution}

\begin{solution}{exo:g11:unemployment:6}
$0.02/0.32 = 6.25\,\%$.
\end{solution}

\begin{solution}{exo:g11:unemployment:7}
$0.015/0.15 = 10\,\%$; each month the gap is multiplied by $1 - 0.15 = 0.85$, so 15\,\% of it closes.
\end{solution}

\begin{solution}{exo:g11:unemployment:8}
$u_1 = 0.1 + 0.3 \times 0.20 = 0.16$; $u_2 = 0.1 + 0.3 \times 0.16 = 0.148$; $u_3 = 0.1 + 0.3 \times 0.148 = 0.1444$. The limit is $0.1/0.7 \approx 14.3\,\%$.
\end{solution}

\begin{solution}{exo:g11:unemployment:9}
$0.02/0.22 \approx 9.1\,\%$, up from 6.25\,\%.
\end{solution}

\begin{solution}{exo:g11:unemployment:10}
United States 24.9\,\% (1933), United Kingdom 15.6\,\% (1932), Germany 17.2\,\% (1932). In 1982: United States 9.7\,\%, United Kingdom 10.9\,\%.
\end{solution}

\begin{solution}{exo:g11:unemployment:11}
$70/200 = 35\,\%$. After twelve months, $0.9^{12} \approx 0.28$: 28\,\% are still unemployed.
\end{solution}

\begin{solution}{exo:g11:unemployment:12}
Unemployment that stays high long after a recession has ended, with no return to its earlier level, as in Britain, at 2.5\,\% in 1970 and 10.9\,\% in 1982, and high for
years after. A natural rate can move when the \omterm{def:g11:unemployment:separation-rate}{job-finding rate} changes for reasons other than the cycle (benefits, the skills of the unemployed, rules on hiring), when the
\omterm{def:g11:unemployment:separation-rate}{separation rate} changes (more short jobs), or when long spells themselves lower the chance of being hired.
\end{solution}

\begin{solution}{exo:g11:unemployment:13}
A shirker gains 30 a month and risks, with probability 0.2, losing $10(w - 1500)$. He does not shirk if $0.2 \times 10 (w - 1500) \ge 30$, that is $w \ge 1515$ dalens.
\end{solution}

\begin{solution}{exo:g11:unemployment:14}
Before: $0.02/0.20 = 10\,\%$. After: $f = 0.18 \times 0.93 \approx 0.167$, and $0.02/0.187 \approx 10.7\,\%$.
\end{solution}

\begin{solution}{exo:g11:unemployment:15}
Young workers change jobs more often and hold more short or trial jobs (a higher \omterm{def:g11:unemployment:separation-rate}{separation rate}); they lack experience and contacts, which lowers their \omterm{def:g11:unemployment:separation-rate}{job-finding rate};
and wage floors or hiring rules weigh most on those whose productivity is still low, again lowering the \omterm{def:g11:unemployment:separation-rate}{job-finding rate}. Differences between countries in how schools hand
young people over to employers (training inside firms, for instance), in the cost of hiring and in \omterm{def:g11:labour-market:minimum-wage}{minimum wages} relative to young people's productivity are candidate
explanations; each needs its own evidence.
\end{solution}

\begin{solution}{pb:g11:unemployment:1}
\textbf{1.} A: $0.01/0.25 = 4\,\%$. B: $0.02/0.20 = 10\,\%$.

\textbf{2.} $0.75 \times 4 + 0.25 \times 10 = 5.5\,\%$.

\textbf{3.} $0.25 \times 10/5.5 \approx 45\,\%$.

\textbf{4.} A: $1/0.24 \approx 4.2$ months; B: $1/0.18 \approx 5.6$ months.

\textbf{5.} With A's \omterm{def:g11:unemployment:separation-rate}{separation rate}, B would have $0.01/0.19 \approx 5.3\,\%$; with A's \omterm{def:g11:unemployment:separation-rate}{job-finding rate}, $0.02/0.26 \approx 7.7\,\%$. The \omterm{def:g11:unemployment:separation-rate}{separation rate} matters more.

\textbf{6.} A: $0.01/0.16 = 6.25\,\%$. B: $0.02/0.12 \approx 16.67\,\%$.

\textbf{7.} After one month: $0.01 + 0.84 \times 0.04 = 4.36\,\%$. After twelve: $6.25 - 2.25 \times 0.84^{12} \approx 6.25 - 2.25 \times 0.123 \approx 5.97\,\%$.

\textbf{8.} $16.67 - 6.67 \times 0.88^{12} \approx 16.67 - 6.67 \times 0.216 \approx 15.23\,\%$.

\textbf{9.} $0.75 \times 5.97 + 0.25 \times 15.23 \approx 8.29\,\%$.

\textbf{10.} From $10 - 4 = 6$ points to $15.23 - 5.97 = 9.26$: by 3.26 points.

\textbf{11.} A: $1 - 0.01 - 0.24 = 0.75$; B: $1 - 0.02 - 0.18 = 0.80$.

\textbf{12.} A: $0.75^n = 0.5$ gives $n = \ln 0.5/\ln 0.75 \approx 2.4$ months; B: $\ln 0.5/\ln 0.8 \approx 3.1$ months.

\textbf{13.} $4 + 1.97 \times 0.75^3 \approx 4 + 0.83 = 4.83\,\%$.

\textbf{14.} Its gap is larger, and it shrinks more slowly each month (by a factor of 0.80 against 0.75), because $s + f$ is smaller.

\textbf{15.} At month 36 the chart's overall rate is 5.51\,\%, against a steady state of 5.5\,\%: the gap is about a hundredth of a point.

\textbf{16.} The excess over the steady state: $8.29 - 5.5 \approx 2.8$ points.

\textbf{17.} $0.85 \times 4 + 0.15 \times 10 = 4.9\,\%$.

\textbf{18.} $0.02/0.18 \approx 11.1\,\%$.

\textbf{19.} The \omterm{def:g11:unemployment:separation-rate}{job-finding rate} falls the longer people are out of work, so the long-term unemployed left by a recession find jobs more slowly than the model assumes; some
also leave the labour force.

\textbf{20.} The steady-state rates are 4\,\% for group A, 10\,\% for group B and 5.5\,\% for Brenland; after the recession, group A closes half its gap in about 2.4 months and
group B in about 3.1.
\end{solution}
